\documentclass[11pt]{article}
\usepackage[margin=1in]{geometry}
\usepackage[T1]{fontenc}\usepackage[utf8]{inputenc}\usepackage{lmodern}
\usepackage{amsmath,amssymb,amsthm,mathtools,booktabs,microtype}
\usepackage[colorlinks=true,linkcolor=blue,urlcolor=blue,citecolor=blue]{hyperref}
\hypersetup{pdftitle={Finite Bayesian neural networks as agnostic approximation learners},pdfauthor={ChatGPT-6 Astra; Claude Opus 5.5; Kaarel Hänni; Dmitry Vaintrob}}
\newtheorem{theorem}{Theorem}[section]\newtheorem{lemma}[theorem]{Lemma}\newtheorem{corollary}[theorem]{Corollary}\newtheorem{proposition}[theorem]{Proposition}
\theoremstyle{remark}\newtheorem{remark}[theorem]{Remark}
\newcommand{\E}{\mathbb E}\newcommand{\R}{\mathbb R}\newcommand{\N}{\mathbb N}\newcommand{\KL}{\operatorname{KL}}\newcommand{\clip}{\operatorname{clip}}\newcommand{\Rad}{\mathfrak R}\newcommand{\op}{\mathrm{op}}
\title{Finite Bayesian neural networks\\as agnostic approximation learners}
\author{ChatGPT-6 Astra \and Claude Opus 5.5 \and Kaarel H\"anni \and Dmitry Vaintrob\thanks{Authors are listed in alphabetical order.}}
\date{17 November 2026}
\begin{document}\maketitle
\begin{abstract}
We prove agnostic squared-loss guarantees for finite Bayesian neural networks with one dense Gaussian architecture and bounded smooth activations, including tanh and sigmoid. The ordinary raw posterior mean competes, with leading coefficient one, against the best smaller bounded-weight network for arbitrary bounded IID labels. A likelihood calibration computed from Bayesian evidence adapts to the comparison network's prediction risk without knowing that risk or the comparison width. With $H\ge2$ hidden layers, let $v_H=\max\{4H,H^2+2H-2\}$ and $\rho_H=2H$. At fixed comparison resources, the excess bound is $\widetilde O(R(g)^{1-\rho_H/v_H}n^{-1/v_H}+n^{-1/\rho_H})$. One sample event covers every learner width above the comparison width and every Gaussian scale at or stronger than mean field. Input dimension enters polynomially. There is no limit in width or prior strength, support restriction, or prior over architectures. The proof controls Gaussian residual fields directly, compresses their retained factors into finite programs, and uses a localized finite-class argument. Separate shallow bounds retain stronger growing-comparison guarantees for monotone activations. Constants depend on depth; the theorem does not give polynomial sample complexity jointly in comparison width and depth.
\end{abstract}

\section{The approximation guarantee}
A large Bayesian neural network should ideally compete with the best smaller network even when the data are not generated by a small network. Such a comparison requires an approximation term, a complexity penalty, and a precise likelihood calibration. We establish this form of guarantee for one fixed fully connected Gaussian architecture at each learner width.

The main technical difficulty near mean-field scaling is that typical internal Gaussian matrices have Frobenius energy proportional to learner width. Their operator norms remain bounded. An argument that charges all Frobenius energy therefore loses width independence, while a direct cover of all Lipschitz functions incurs a curse of input dimension. We instead control the Gaussian residual feature operator directly and compress retained features in a Hilbert space of functions. The Hilbert-space measure may combine the population input law with the empirical input distribution; a fixed family of proof programs supplies the statistical event.

Vaintrob and H\"anni~\cite{vh} study width-robust mean-field learnability through reduced entropy and network approximation. The theorem here specifies an agnostic squared-loss oracle for every finite width and every scale at or stronger than mean field, on arbitrary input distributions in a continuous domain, with explicit depth-dependent rate exponents. Its likelihood calibration and raw-output prediction are stated below.

Fix a twice continuously differentiable activation $\phi:\R\to[-1,1]$ with $|\phi'|\le1$ and bounded second derivative. This includes tanh and logistic sigmoid. We do not assume $\phi(0)=0$. Let $H$ denote the number of hidden layers and $L=H+1$ the number of affine maps, including the linear readout. The main theorem permits every hidden depth $H\ge1$. Inputs lie in $[-1,1]^d$, $d\ge1$, and labels in $[-1,1]$. Their joint law $D$ is otherwise arbitrary, and the training sample $S=((X_i,Y_i))_{i=1}^n$ is IID from $D$. For a real prediction $f$, write
\[
 R(f)=\E_D(f(X)-Y)^2,\qquad \widehat R_S(f)=n^{-1}\sum_i(f(X_i)-Y_i)^2.
\]
A hidden-width-$N$ network has
\[
 h^1(x)=\phi(W^1\bar x),\quad
 h^\ell(x)=\phi(W^\ell h^{\ell-1}(x))\ (2\le\ell\le H),\quad
 z_w(x)=v^\top h^H(x),\qquad \bar x=(1,x).
\]
There are no internal biases. A teacher uses the same depth and activation, width $m$, and coefficients bounded by $B\ge1$. For tanh, teachers with internal and readout biases are included by one additional constant neuron per hidden layer and a coefficient cap $2B$: carry the constant $1/2$, whose copying weight is $2\operatorname{arctanh}(1/2)=\log3$, and replace each bias $b$ by an edge of weight $2b$.

For $0<\tau\le1$, the learner prior $\mu_{N,\tau}$ has independent centered Gaussian coordinates with variances
\begin{equation}
 \tau^2/d\quad\text{in }W^1,\qquad
 \tau^2/N\quad\text{in }W^2,\ldots,W^H,\qquad
 \tau^2/N^2\quad\text{in }v.
 \label{eq:prior}
\end{equation}
The likelihood is Gaussian with observation variance $n\lambda\tau^2/N$. Thus the actual finite posterior is
\begin{equation}
 \Pi_{N,\tau,S}(dw)=Z^{-1}
    \exp\!\left[-\frac{N\widehat R_S(z_w)}{2\lambda\tau^2}\right]
    \mu_{N,\tau}(dw).
 \label{eq:posterior}
\end{equation}
The canonical exponent $\gamma\ge1$ is obtained by setting $\tau^2=N^{1-\gamma}$. Mean field is $\tau=1$. The prediction is the ordinary posterior mean $\bar z_{N,\tau,S}=\int z_w\,d\Pi$.

For a teacher $g$, define
\begin{gather}
 E_g=d(d+1)B^2+2Hm^2B^2,\qquad U_g=mB,\qquad D_g=d+2+(H-1)m,\notag\\
 b_0=2\sqrt{2(2\log9+\log(4(H+1)))},\qquad
 P_g=1+2mB(d+H+2)(2+3mB+b_0)^{H-1},\notag\\
 \rho=\frac{\lambda}{4(U_g+1)},\qquad
 \eta_g=\min\{1/2,\rho/(2P_g)\},\qquad
 A_g(\lambda)=\frac{E_g}{2}+D_g\log\frac{8}{3\eta_g}+\log2.
 \label{eq:teacher-cost}
\end{gather}
Constants written $C_H,c_H>0$ depend only on $H$ and a bound on $|\phi''|$, and can change between displays. For $H\ge2$, put
\[
 p_H=H^2+2H-3,\quad q_H=4H-2,\quad
 \rho_H=2H,\quad v_H=\max\{4H,H^2+2H-2\},\quad r_H=v_H/(2H).
\]
At $H=1$, set $v_1=2$, $\rho_1=1$, and $r_1=2$. Thus $r_H\ge2$ at every depth. Define the empirical free energy
\[
 F_{N,\tau,S}(\lambda)=-\frac{2\lambda}{\omega}
 \log\int e^{-\omega\widehat R_S(z_w)/(2\lambda)}\,d\mu_{N,\tau}(w),
 \qquad \omega=N/\tau^2.
\]
All quantities needed for the following calibration depend only on the specified learner, sample, and confidence. Fix sufficiently large depth-dependent constants and set
\[
 b_n=C_H\log(C_Hn/\delta),\qquad
 \theta_n=\frac{C_H(d+1)^2[\log(C_Hn(d+1)/\delta)]^{12H+30}}n.
\]
When $\theta_n b_n^{v_H}5^{v_H-\rho_H}\le1$, choose the unique positive solution
\begin{equation}
 \lambda_{N,\tau,S}^{v_H}
 =\theta_n[2F_{N,\tau,S}(\lambda_{N,\tau,S})+
                  b_n\lambda_{N,\tau,S}]^{v_H-\rho_H}.
 \label{eq:main-calibration}
\end{equation}
Otherwise set $\lambda_{N,\tau,S}=2\omega$. This fallback only concerns sample sizes below a depth-dependent threshold. The likelihood variance is $n\lambda_{N,\tau,S}\tau^2/N$ in both cases. Section~\ref{sec:localized-direct-oracle} proves uniqueness, measurability, and the risk guarantee of this calibration.

\begin{theorem}[An evidence-calibrated oracle at every finite width]
\label{thm:main}
Let $H\ge1$, $n\ge2$, and $0<\delta<1$. For each fixed width-$m$, coefficient-cap-$B$ teacher $g$, set
\[
 \mathcal A_{g,n}=C_H+2A_g(1/n)+\log(C_Hn/\delta),
\]
with $A_g$ from \eqref{eq:teacher-cost}. With probability at least $1-\delta$, simultaneously for every learner width $N\ge m$ and every $0<\tau\le1$, the posterior with calibration \eqref{eq:main-calibration} satisfies
\begin{align}
 \int R(z_w)\,d\Pi_{N,\tau,S}
 \le R(g)+C_H\big[&\mathcal A_{g,n}
           R(g)^{1-1/r_H}\theta_n^{1/v_H}
       +\mathcal A_{g,n}^{r_H}\theta_n^{1/\rho_H}\big]+n^{-3}.
 \label{eq:main-oracle}
\end{align}
The same inequality holds for the squared risk of the ordinary raw posterior mean.
\end{theorem}

The learner does not know $g$, its resource bounds, the input law, or $R(g)$. A population-risk minimizer in a compact bounded-weight teacher class can be substituted for $g$, preserving the leading coefficient one. Countably many declared comparison classes can be covered on one event by allocating confidence. No realizability or observation-noise-model assumption is made. Outputs and the prior remain unrestricted; clipping occurs only in proof approximants.

\begin{corollary}[Rates and growing comparison networks]
\label{cor:rates}
At fixed teacher resources, dimension, depth, and confidence, the excess is
\[
 \widetilde O_H\!\left(R(g)^{1-\rho_H/v_H}n^{-1/v_H}+n^{-1/\rho_H}\right).
\]
Thus a fixed positive comparison risk gives rate $n^{-1/v_H}$, while a zero-risk comparison gives rate $n^{-1/\rho_H}$, up to logarithms. At fixed $d,H$, comparison width $n^\alpha$ and coefficient cap $n^\beta$ are allowed with vanishing excess whenever $\alpha+\beta<1/(2v_H)$.
\end{corollary}

\begin{center}
\begin{tabular}{ccc}
\toprule
Hidden layers $H$ & Fixed positive comparison risk & Zero comparison risk\\
\midrule
1 & $\widetilde O(n^{-1/2})$ & $\widetilde O(n^{-1})$\\
2 & $\widetilde O(n^{-1/8})$ & $\widetilde O(n^{-1/4})$\\
3 & $\widetilde O(n^{-1/13})$ & $\widetilde O(n^{-1/6})$\\
4 & $\widetilde O(n^{-1/22})$ & $\widetilde O(n^{-1/8})$\\
5 & $\widetilde O(n^{-1/33})$ & $\widetilde O(n^{-1/10})$\\
\bottomrule
\end{tabular}
\end{center}
These are sufficient exponents, not optimality claims. Section~\ref{sec:one-hidden} also gives fixed-calibration bounds; monotone activations permit the larger growing-resource range $\alpha+\beta<1/2$.

\section{A comparison distribution around a small network}
\label{sec:teacher}
Write $\omega=N/\tau^2$.  The quadratic form in the Gaussian density is
\[
 Q_N(w)=\frac dN\|W^1\|_F^2+\sum_{\ell=2}^H\|W^\ell\|_F^2+N\|v\|_2^2,
 \qquad d\mu_{N,\tau}\propto e^{-\omega Q_N/2}\,dw.
\]
Partition the $N$ neurons in each hidden layer into $m$ groups of sizes $k_j\in\{\lfloor N/m\rfloor,\lceil N/m\rceil\}$. Copy the teacher's first-layer rows, divide every internal edge from group $j$ by $k_j$, and divide readout coefficients similarly. This balanced clone computes $g$ exactly. The inequality $N/k_j\le2m$ gives $Q_N(w_*)\le E_g$.

\begin{lemma}[Teacher cylinder]
\label{lem:teacher}
For the quantities in \eqref{eq:teacher-cost}, the prior assigns mass at least $e^{-\omega A_g(\lambda)}$ to the event
\[
 U_\rho(g)=\{w:\sup_{x\in[-1,1]^d}|z_w(x)-g(x)|<\rho\}.
\]
Consequently, a comparison law $\nu_g$ satisfies
\begin{equation}
 \int\widehat R_S(z_w)d\nu_g+\frac{2\lambda}{\omega}\KL(\nu_g\Vert\mu_{N,\tau})
 \le\widehat R_S(g)+\lambda+2\lambda A_g(\lambda).
 \label{eq:comparison}
\end{equation}
\end{lemma}
\begin{proof}
Use standard Gaussian coordinates for every parameter block. Let $P$ project onto the group-constant subspace, with orthonormal basis $u_j=k_j^{-1/2}\mathbf1_{\text{group }j}$. Impose absolute value at most $\eta_g$ on every first-layer and readout standardized coordinate and on every internal row coefficient $\langle Z_{i,*},u_j\rangle$. These $ND_g$ independent constraints have probability at least $(3\eta_g/8)^{ND_g}$. Indeed a standard Gaussian has mass at least $3\eta/8$ on $[-\eta,\eta]$ for $0<\eta\le1$, by unimodality and its mass at least $3/4$ on $[-2,2]$.

Conditioning a standard Gaussian coordinate to a symmetric interval preserves the bound $\E e^{sZ}\le e^{s^2/2}$; this follows either by differentiating its truncated moment-generating function or by pairing its two symmetric tails. Orthogonal coordinates factor, so all centered linear forms remain unit-subgaussian. Two $1/4$-nets of the Euclidean sphere, each of size at most $9^N$, imply
\[
 \Pr(\|Z\|_\op>2u\sqrt N\mid\text{cylinder})
 \le2\exp[(2\log9-u^2/2)N].
\]
Intersect with $\|Z\|_\op\le b_0\sqrt N$ at every internal layer. A union bound leaves a centrally symmetric event $A$ with mass at least $\tfrac12(3\eta_g/8)^{ND_g}$.

Translate this event by the standardized clone. Internal perturbations are $\tau Z/\sqrt N$. Their part on $I-P$ annihilates the clone's hidden states, and their part on $P$ has operator norm at most $\eta_g\sqrt m$. With normalized hidden-state errors $e_\ell=\|h^\ell-h_*^\ell\|_2/\sqrt N$, the Lipschitz property of $\phi$ gives
\[
 e_1\le\eta_g(d+1),\qquad
 e_\ell\le(\sqrt2mB+b_0)e_{\ell-1}+\eta_g\sqrt m.
\]
The readout perturbation has $\ell^1$ norm at most $\eta_g$, and $\sqrt N\|v_*\|_2\le\sqrt2mB$. The resulting uniform output error is at most $\eta_g P_g\le\rho/2$.

For a standard Gaussian measure $G$, symmetry gives
\[
 G(A+a)\ge G(A)e^{-\|a\|^2/2},
\]
since the density ratio is $e^{-\|a\|^2/2}e^{-\langle Z,a\rangle}$ and its conditional expectation on $A$ is at least its value at the conditional mean zero. Here $\|a\|^2=\omega Q_N(w_*)\le\omega E_g$. As $N\le\omega$ and $1\le\omega$, the resulting mass bound is the claimed one.

Take $\nu_g=\mu_{N,\tau}(\cdot\mid U_\rho(g))$. Its entropy is at most $\omega A_g$, and its empirical risk is at most $\widehat R_S(g)+2\rho(U_g+1)+\rho^2\le\widehat R_S(g)+\lambda$. This proves \eqref{eq:comparison}.
\end{proof}

\section{A Gaussian residual operator bound}
\label{sec:residual}
Fix any probability measure $\xi$ on the input cube. It need not be known to the learner. All Hilbert norms in this section are in $L^2(\xi)$.

For an actual hidden layer, let
\[
 Q(x,y)=N^{-1}\sum_{i=1}^Nh_i(x)h_i(y),
 \qquad T c=N^{-1/2}\sum_i c_i h_i.
\]
Then $Q=TT^*$ as an integral operator and $\operatorname{tr}Q\le1$. At a spectral cutoff $\eta$, retain eigenvalues at least $\eta$, with factors $F_j=\sqrt{\lambda_j}e_j$. There are at most $1/\eta$ of them. The corresponding right-singular vectors $u_j\in\R^N$ are orthonormal and satisfy $T u_j=F_j$. Thus
\[
 F_j=N^{-1}\sum_i b_{ij}h_i,\qquad b_{ij}=\sqrt N u_j(i),
 \qquad N^{-1}\sum_i b_ib_i^\top=I.
\]
The residual kernel $E=Q-\sum_jF_j\otimes F_j$ is positive, has operator norm at most $\eta$, and has pointwise variance
$v(x)=Q(x,x)-\|F(x)\|^2\in[0,1]$.

Conditional on this preceding layer, the next Gaussian preactivation rows have the decomposition
\begin{equation}
 Z_i(x)=a_i\cdot F(x)+Z_i^0(x),
 \quad a_i\sim N(0,\tau^2I),\quad Z_i^0\sim\mathrm{GP}(0,\tau^2E),
 \label{eq:direct-decomposition}
\end{equation}
independently over rows and between the two parts. Define
\[
 m_i(x)=\E[\phi(Z_i(x))\mid a_i],\qquad
 r_i(x)=\phi(Z_i(x))-m_i(x),\qquad
 R c=N^{-1/2}\sum_i c_i r_i.
\]

\begin{lemma}[Residual feature concentration]
\label{lem:direct-residual}
For $K\ge1$, conditional on the preceding layer and every value of the latent vectors $(a_i)$, with probability at least $1-2e^{-NK/\tau^2}$,
\begin{align}
 \|R\|_\op&\le C[\omega^{-1/2}+\sqrt{\eta(K+1)}],\notag\\
 \left\|N^{-1}\sum_i\left(\phi(Z_i)^2-\E[\phi(Z_i)^2\mid a_i]\right)\right\|_{L^2(\xi)}
 &\le C[\omega^{-1/2}+\sqrt{\eta K}].
 \label{eq:direct-residual}
\end{align}
\end{lemma}
\begin{proof}
Fix a unit vector $c\in\R^N$. Conditional residual rows are independent and centered, and Gaussian Poincar\'e gives conditional pointwise variance at most $\tau^2v(x)$ because $\phi$ is 1-Lipschitz. Thus each residual row has integrated second moment at most $\tau^2$, and $\E\|Rc\|\le\tau/\sqrt N=\omega^{-1/2}$. As a function of all residual Gaussian coordinates, $\|Rc\|$ is Lipschitz with constant at most $\tau\sqrt{\eta/N}$. To verify this, dualize the Hilbert norm. The derivative of $\langle f,Rc\rangle$ in row $i$ has squared Gaussian-coordinate norm
\[
 \frac{c_i^2\tau^2}{N}
 \langle f\phi'(Z_i),E(f\phi'(Z_i))\rangle
 \le\frac{c_i^2\tau^2\eta}{N}\|f\|_2^2.
\]
Sum over rows and take $\|f\|_2=1$. Gaussian concentration therefore gives
\[
 \Pr(\|Rc\|>\omega^{-1/2}+t\mid(a_i))
 \le e^{-Nt^2/(2\tau^2\eta)}.
\]
A $1/2$-net of the unit sphere of $\R^N$ has at most $5^N$ points, and the operator norm is at most twice the maximum over this net. Choose $t=\sqrt{2\eta(K+\tau^2\log5)}$ and use a union bound. This proves the first assertion, uniformly over all coefficient vectors, including coefficients subsequently selected from the same residual realization.

For the second, Gaussian Poincar\'e gives the centered squared-feature Hilbert average expected norm at most $2\omega^{-1/2}$. Since $\phi^2$ is 2-Lipschitz, its norm is Gaussian-Lipschitz with constant at most $2\tau\sqrt{\eta/N}$. Concentration gives the stated bound directly. Infinite-dimensional Gaussian fields cause no additional issue: here every residual field is a linear image of finitely many row-weight coordinates, and the Hilbert-valued maps are Lipschitz.
\end{proof}

The standard Gaussian norm event
\begin{equation}
 \|W^1\|_F^2\le NK,\qquad
 \max_{2\le\ell\le H}\|W^\ell\|_\op^2\le K,\qquad
 N\|v\|^2\le K
 \label{eq:norm-event}
\end{equation}
has prior failure at most $(H+1)e^{-NK/(16\tau^2)}$ for a sufficiently large numerical lower bound on $K$. Its proof uses a chi-square tail for the first layer and readout and the Gaussian operator-norm tail for each internal matrix.

There is also an exact latent Gram bound on this event. If $V$ has the retained preceding right-singular vectors as columns, then $a_i=\sqrt N(WV)_{i,*}$, so
\begin{equation}
 N^{-1}\sum_i a_ia_i^\top=V^\top W^\top WV\preceq K I.
 \label{eq:actual-latent-gram}
\end{equation}
No independence between the retained current factors and their residuals is assumed. The sphere net in Lemma~\ref{lem:direct-residual} was taken precisely to control these adaptively selected coefficient vectors.

\section{Finite programs approximating the good networks}
\label{sec:compression}
Assume $H\ge2$. The programs below only bound statistical capacity; they are not the Bayesian learner. Define
\[
 \psi(m,v)=\E_G\phi(m+\sqrt vG),\qquad
 \chi(m,v)=\E_G\phi(m+\sqrt vG)^2,\quad G\sim N(0,1).
\]
The Gaussian heat equation shows that these functions are Lipschitz in both arguments, including at $v=0$, with constants depending only on $\sup|\phi''|$. They have ranges $[-1,1]$ and $[0,1]$, respectively, and $\chi\ge\psi^2$.

\begin{proposition}[One fixed program family for every width and scale]
\label{prop:programs}
For $K\ge K_H$ and $0<a\le1$, there is a compact finite-parameter family $\mathcal P(K,a)$ of functions bounded by $\sqrt K$, independent of the input measure, width, scale, and training sample, whose scalar parameter count is at most
\begin{equation}
 P_H(K,a)=C_H(d+1)K^{p_H}a^{-q_H},\qquad
 p_H=H^2+2H-3,\quad q_H=4H-2.
 \label{eq:program-count}
\end{equation}
Its uniform covering entropy obeys
\begin{equation}
 \log\mathcal N(u,\mathcal P(K,a),\|\cdot\|_\infty)
 \le C_HP_H(K,a)\log\frac{C_H(d+1)K}{au}
 \quad(0<u\le\sqrt K).
 \label{eq:program-entropy}
\end{equation}
For every fixed probability measure $\xi$ and every $N\ge1$, $0<\tau\le1$, there is a Borel set $G_{N,\tau,\xi}(K,a)$ such that
\begin{equation}
 \mu_{N,\tau}(G^c)\le C_He^{-\omega K/32},\qquad \omega=N/\tau^2,
 \label{eq:direct-prior-good}
\end{equation}
and every output in this set is bounded by $\sqrt K$ and has an approximant $f\in\mathcal P(K,a)$ with $\|z_w-f\|_{L^2(\xi)}\le C_Ha$.
\end{proposition}

\subsection{Separate factor and variance approximations}
Retain factors $F_\ell$ at cutoff $\eta_\ell$, with rank $r_\ell\le\eta_\ell^{-1}$, and write $v_\ell=Q_\ell(x,x)-\|F_\ell(x)\|^2$. Actual factors have pointwise norm at most one, and $v_\ell\in[0,1]$. Use finite averages over actual rows, denoted $\E_i$, and exact coefficients satisfying $\E_i b_i b_i^\top=I$.

For any preceding factor function $F$, including an unbounded approximant, and any $[0,1]$-valued variance function $v$, set
\[
 F_\ell^0=\E_i b_i\psi(a_i\cdot F,\tau^2v),\qquad
 v_\ell^0=\E_i\chi(a_i\cdot F,\tau^2v)-\|F_\ell^0\|^2.
\]
Bessel's inequality and $\chi\ge\psi^2$ give
\[
 \|F_\ell^0(x)\|^2\le\E_i\psi(a_i\cdot F(x),\tau^2v(x))^2
 \le\E_i\chi(a_i\cdot F(x),\tau^2v(x))\le1.
\]
Hence $v_\ell^0\in[0,1]$. For two preceding states, the vector Bessel inequality and \eqref{eq:actual-latent-gram} bound the changes in both $F_\ell^0$ and $v_\ell^0$ by
\begin{equation}
 C_H\big[\sqrt K\|\Delta F\|_{L^2(\xi;\R^{r_{\ell-1}})}+\|\Delta v\|_{L^2(\xi)}\big].
 \label{eq:separate-propagation}
\end{equation}
For the variance use $|\|F_\ell^0\|^2-\|\widetilde F_\ell^0\|^2|\le2\|F_\ell^0-\widetilde F_\ell^0\|$. Both exact averaged vectors are pointwise bounded, even if the preceding sampled vector is not.

Replacing actual current features by their conditional means changes the factor vector by at most
\[
 E_\ell=C_H\sqrt{r_\ell+1}\,[\omega^{-1/2}+\sqrt{K\eta_{\ell-1}}]
\]
on the residual event of Lemma~\ref{lem:direct-residual}. The diagonal changes by this bound without the rank factor. Thus the actual variance differs from $v_\ell^0$ by at most $C_HE_\ell$ as well.

Evaluate these exact means on already chosen preceding approximants. Approximate $F_\ell^0$ by sampling row indices with replacement. Since $\E_i\|b_i\psi_i\|^2\le r_\ell$, a sample of
\[
 m_\ell^F=C_H(r_\ell+1)/(\delta_\ell^F)^2
\]
suffices for vector error $\delta_\ell^F$. Leave this vector approximation linear: do not project it onto a unit ball.

The variance admits a separate scalar representation
\begin{equation}
 v_\ell^0=\E_i\chi_i-\E_{i,j}(b_i\cdot b_j)\psi_i\psi_j,
 \qquad \E_{i,j}(b_i\cdot b_j)^2=r_\ell,
 \label{eq:pair-variance}
\end{equation}
where $i,j$ are independent uniform row indices. The identity follows by taking the trace of $(\E_i b_i b_i^\top)^2$. Independent samples of each scalar mean, with total size
$m_\ell^V=C_H(r_\ell+1)/(\delta_\ell^V)^2$, approximate it to error $\delta_\ell^V$. Store each pair coefficient $b_i\cdot b_j$ as one scalar and clip only this variance approximation to $[0,1]$.

These positive-probability sampling assertions follow from the Hilbert-space second-moment bound and Markov's inequality. The second moments of $a,b,b_i\cdot b_j$ are at most $Kr_{\ell-1},r_\ell,r_\ell$, respectively. Additional Markov bounds on all selected coefficient sums provide samples with polynomially bounded scalar parameters while preserving positive probability of the approximation event. Sampling is performed after preceding approximants are fixed, so no empirical Gram claim is required.

At the first layer apply the same construction directly to the finite features $h_i=\phi(w_i\cdot\bar x)$, using
$F_1=\E_i b_i h_i$ and $v_1=\E_i h_i^2-\E_{i,j}(b_i\cdot b_j)h_ih_j$.
The first-layer norm event gives $\E_i\|w_i\|^2\le K$. There is no residual deletion and no minimum-width condition at this layer: sampling with replacement is allowed.

Writing $e_\ell^F,e_\ell^V$ for the two errors, for $\ell\ge2$ the construction gives
\begin{align*}
 e_\ell^F&\le C_H(\sqrt K e_{\ell-1}^F+e_{\ell-1}^V+E_\ell+\delta_\ell^F),\\
 e_\ell^V&\le C_H(\sqrt K e_{\ell-1}^F+e_{\ell-1}^V+E_\ell+\delta_\ell^V).
\end{align*}
At layer one only the respective sampling errors occur.

\subsection{Terminal readout and effective-width cutoff}
The actual output is $z=\E_i c_i h_i^H$, with $c_i=Nv_i$ and $\E_i c_i^2\le K$. Its residual deletion error is at most
$C_H[\sqrt{K/\omega}+K\sqrt{\eta_{H-1}}]$.
Propagation through its conditional mean costs at most
$C_H[K e_{H-1}^F+\sqrt K e_{H-1}^V]$.
A sample of $m_H=C_HK/a^2$ final rows then gives scalar error $c_Ha$. Clip the final output to $[-\sqrt K,\sqrt K]$.

For $t=H-\ell$ choose sufficiently small depth-dependent constants in
\[
 \delta_\ell^F\asymp_H aK^{-(t+1)/2},\qquad
 \delta_\ell^V\asymp_H aK^{-t/2},\qquad
 \eta_\ell\asymp_H a^{2t}K^{-t(t+3)/2}.
\]
The recurrence closes because $\eta_{\ell-1}\le c_H(\delta_\ell^F)^2\eta_\ell/K$. All residual sampling terms are also small enough if
\begin{equation}
 \omega\ge T_H(K,a):=C_HK^{s_H}a^{-(2H-2)},\qquad
 s_H=(H^2+H-4)/2.
 \label{eq:direct-width}
\end{equation}
Indeed the largest intermediate requirement $(r_\ell+1)/(\delta_\ell^F)^2$ starts at $\ell=2$; include the final requirement $K/a^2$. At $H=2$ the latter is the only one. The output approximation error is then at most $C_Ha$.

\subsection{Contracting factors into subsequent coefficients}
The retained dimensions need not appear among the program's parameters. The sampled vector $\widetilde F_\ell$ is a linear combination of $m_\ell^F$ scalar sampled features. Every occurrence of this vector is a linear form $a\cdot\widetilde F_\ell$ in the next layer. Contract $a$ with its sampled vector coefficients and store the resulting $m_\ell^F$ scalar coefficients. This is possible because no radial projection is applied. The separate variance formula \eqref{eq:pair-variance} also avoids evaluating the norm of the sampled vector.

The total number of scalar features at layer $\ell$ is at most $m_\ell^F+3m_\ell^V\le C_Hm_\ell^F$. Consequently the parameter count is at most
\[
 C_H\left[(d+1)m_1^F+\sum_{\ell=2}^H m_\ell^F m_{\ell-1}^F\right],\qquad m_H^F=m_H.
\]
For intermediate $t=H-\ell$,
\[
 m_\ell^F\le C_HK^{(t^2+5t+2)/2}a^{-2t-2},\qquad
 m_\ell^V\le C_HK^{(t^2+5t)/2}a^{-2t-2}.
\]
The largest adjacent product, $m_1^Fm_2^F$, gives exactly \eqref{eq:program-count}; when $H=2$ it is $K^5a^{-6}$. All contracted coefficients remain polynomially bounded. Gaussian mean features are bounded, and variances are clipped to $[0,1]$, so parameter-to-function Lipschitz constants are polynomial in $d+1,K,a^{-1}$ at fixed depth. A uniform parameter grid proves \eqref{eq:program-entropy}.

Intersecting the norm event with each conditional residual event proves \eqref{eq:direct-prior-good} when \eqref{eq:direct-width} holds. Measurable spectral bases can be chosen from finite Gram matrices, including repeated eigenvalues, so the tests are Borel.

For the remaining effective widths $\omega<T_H(K,a)$, necessarily $N<T_H(K,a)$. Include in the same program family the actual finite networks of every integer width below this cutoff on the norm event. Their weights have polynomial bounds, and their parameter count is at most
\[
 C_H[dT_H+T_H^2]\le C_H(d+1)K^{p_H}a^{-q_H},
\]
because $2s_H=H^2+H-4<p_H$ and $4H-4<q_H$. Their parameter-to-function Lipschitz constants are polynomial as well. A finite union changes only the constant in the entropy bound. This proves the proposition at every finite width and scale, without a second statistical argument.

\section{A localized oracle from Gaussian compression}
\label{sec:localized-direct-oracle}
This section is an abstract statistical consequence of the direct Gaussian
compression lemma. It keeps the same fixed architecture and prior.
Let $\omega\ge1$ be the precision, and let
\[
 d\Pi_{S,\lambda}=Z^{-1}e^{-\omega\widehat R_S(z)/(2\lambda)}d\mu,
 \qquad F_S(\lambda)=-2\lambda\log Z/\omega.
\]
Inputs have any fixed probability law, and labels satisfy $|Y|\le1$.
Suppose uniformly over all eligible widths and scales:
\begin{enumerate}
\item $\mathbb E_\mu\widehat R_S\le M$ for every sample, with a fixed
constant $M\ge1$.
\item For every $K\ge K_0$, $0<a\le1$, and probability measure $\xi$
on inputs, a measurable set $G_\xi(K,a)$ has
$\mu(G_\xi(K,a)^c)\le C e^{-c\omega K}$. Every member of this set has
amplitude at most $C K^{b/2}$ and an $L^2(\xi)$ approximant, with error
at most $a$, in a fixed finite dictionary $\mathcal D(K,a)$ independent
of $\xi$, width, and scale. The dictionary has amplitude at most
$C K^{b/2}$ and
\[
 \log|\mathcal D(K,a)|\le B K^p a^{-q}
       \left[\log\frac{C K}{a}\right]^s.
\]
\item These good sets and approximants are measurable when
$\xi=(D_X+\widehat D_{X,S})/2$. Constants may depend on the fixed depth
and activation, while $B$ is polynomial in input dimension.
\end{enumerate}
Assume $q\ge0$ and enlarge $B$ so that $B\ge1$.
Here $b>0$ is the amplitude exponent. Put $\rho=(q+2)/2$,
$v=\max\{p+b,2\rho\}$, and $r=v/\rho\ge2$.
For the general Lipschitz construction, $b=L$ and $p=Lq/2$, so
$v=L\rho$. Bounded activations allow $b=1$.
All assertions below permit a countable union of fixed teacher classes
by summable confidence allocation.

\subsection{One sample event for all calibrations}
Set $U_0=2M+1$. There are constants $B_0,B_1,b_0,C_0$ determined only
by the assumptions, such that, with
\[
 b_n=b_0+4\log(C_0 n^3),\qquad
 \theta_n=\frac{C_0 B[\log(C_0 n/\delta)]^{s+4}}{n},
\]
the following learner is well defined. If
$\theta_n b_n^v U_0^{v-\rho}\le1$, choose the positive solution of
\begin{equation}
 \lambda_S^v=\theta_n[2F_S(\lambda_S)+b_n\lambda_S]^{v-\rho}.
 \label{eq:localized-evidence-root}
\end{equation}
Otherwise choose $\lambda_S=\omega M$.
The root is unique and measurable because
$\lambda^v/[2F_S(\lambda)+b_n\lambda]^{v-\rho}
=\lambda^\rho/[2F_S(\lambda)/\lambda+b_n]^{v-\rho}$
is continuous and strictly increasing from zero to infinity.
Here $F_S\le M$ and $F_S(\lambda)/\lambda$ is nonincreasing.
On the root branch, enlarging $C_0,b_0$ if necessary gives
$1/n\le\lambda_S\le1/b_n$ and
$u_S:=2F_S(\lambda_S)+b_n\lambda_S\le U_0$.

For use in the proof, prebuild dyadic norm levels $K_j=K_0 2^j$
through $C_0n$ and dyadic accuracies from $c_0/n$ through one.
There is a single sample event of probability at least $1-\delta/2$
with the following property. Whenever $\lambda\ge1/n$,
$0<u\le U_0$, $\lambda K\le C_1u$, and $a$ lies within a factor two
of $c_1\lambda K/\sqrt u$, each network good at $(K,a)$ with
$\widehat R_S(z)\le u$ satisfies
\begin{equation}
 R(z)\le\widehat R_S(z)+\beta\lambda K
 \label{eq:localized-level-deviation}
\end{equation}
provided
\begin{equation}
 n\ge C B[\log(Cn/\delta)]^{s+4}
       u^\rho\lambda^{-2\rho}K^{p+b-2\rho}.
 \label{eq:localized-level-budget}
\end{equation}
Here $\beta>0$ is a sufficiently small fixed constant, and $c_1$
is chosen after $\beta,C_1,U_0$, so all required accuracies are at most
one and at least $c_0/n$.

Indeed, apply finite-class Bernstein to the squared losses of every
fixed dictionary, allocating confidence over the finite level/accuracy
list. A dictionary member has loss range $O(K^b)$ and variance at most
that range times its population risk. Inverting Bernstein gives
\[
 R(f)\le\widehat R(f)+C\sqrt{\widehat R(f)K^b t/n}+C K^b t/n,
 \quad t=\log|\mathcal D|+\log(1/\delta_{j,a}).
\]
The $L^2((D_X+\widehat D_X)/2)$ approximation controls both empirical
and population root risks. At the selected accuracy the approximant
has empirical risk $O(u)$, and transferring root risks costs
$O(\sqrt u\,a+a^2)=O(c_1\lambda K)$.
The leading Bernstein condition is exactly
\eqref{eq:localized-level-budget}; its linear term is absorbed using
$\lambda K\le C_1u$. Constants can be chosen to yield
\eqref{eq:localized-level-deviation}.
If $p+b\ge2\rho$, maximize the right side of
\eqref{eq:localized-level-budget} at $K\le C_1u/\lambda$.
If $p+b<2\rho$, maximize it at the fixed lower cutoff $K_0$.
In both cases the resulting sufficient budget is
$C B\log^{s+4}(Cn/\delta)u^{v-\rho}\lambda^{-v}$.
Thus the evidence root enforces all required level budgets.

\subsection{A sample-dependent complexity with uniform prior tails}
Fix a sample on the preceding statistical event and its selected root
$\lambda,u$. For every dyadic level $K_j$ whose value can enter
$K_j\le u/(\beta\lambda)$, choose the prebuilt accuracy nearest below
$c_1\lambda K_j/\sqrt u$. At larger levels set the accuracy to one.
Use the good sets with $\xi_S=(D_X+\widehat D_X)/2$ and these accuracies.
Define
\[
 Q_S(w)=\max\{K_0,\ 2\sup\{K_j:w\notin G_{\xi_S}(K_j,a_j)\}\},
\]
where the supremum of the empty set is zero. For each fixed sample the
prior tail bound and a countable union give
\[
 \mu(Q_S>t)\le C'e^{-c'\omega t}\quad(t\ge2K_0).
\]
Only finitely many levels fail almost surely, so $Q_S$ is finite
almost surely. If $Q_S=K_j$, then $w\in G_{\xi_S}(K_j,a_j)$.
In particular $R(z_w)\le C(1+Q_S^b)$ almost surely.
Choose $\beta$ small enough that tail integration gives, pointwise
for every sample,
\begin{equation}
 \log\mathbb E_\mu e^{\beta\omega Q_S}\le\omega B_0,
 \qquad
 \mathbb E_\mu[R(z_w)e^{\beta\omega Q_S/4}]
       \le C_0e^{\omega B_1}.
 \label{eq:localized-complexity-mgf}
\end{equation}
The constants do not depend on the sample, selected root, width, or
scale. Polynomial factors in $Q_S$ are absorbed by a smaller
exponential rate, using $\omega\ge1$.
Measurability follows from the measurable good-set tests and countable
supremum. No data-dependent complexity is placed inside an
unjustified empirical-process supremum; its prior inequalities are
pointwise in the sample.

Let $J=\widehat R_S+\beta\lambda Q_S$. On $J>u$,
\[
 e^{-\omega\widehat R_S/(2\lambda)}
 \le e^{-\omega u/(4\lambda)}e^{\beta\omega Q_S/4}.
\]
After normalization and \eqref{eq:localized-complexity-mgf}, the
posterior risk on this set is at most
$C_0\exp[\omega F_S/(2\lambda)-\omega u/(4\lambda)+\omega B_1]$,
which is at most $n^{-3}$ by the choice of $b_n$.
On its complement, \eqref{eq:localized-level-deviation} applies at
$K=Q_S$ and gives $R\le\widehat R_S+\beta\lambda Q_S$.
Entropy duality in \eqref{eq:localized-complexity-mgf} gives
\[
 \beta\lambda\mathbb E_\Pi Q_S
 \le(\lambda/\omega)\operatorname{KL}(\Pi\Vert\mu)+B_0\lambda.
\]
Integrate, extend the nonnegative right side to the entire posterior,
and use the variational identity. This proves the observable certificate
\begin{equation}
 \int R(z_w)d\Pi_{S,\lambda_S}
 \le F_S(\lambda_S)+B_0\lambda_S+n^{-3}.
 \label{eq:localized-certificate}
\end{equation}
The same holds for the ordinary raw posterior mean by Jensen. On the
fallback branch, Jensen gives $Z\ge e^{-1/2}$, so the deterministic
certificate $e^{1/2}M$ is valid instead.

\subsection{The resulting oracle powers}
Suppose a fixed teacher $g$ has a deterministic comparison resource
$A\ge1$ with
$F_S(\lambda)\le\widehat R_S(g)+\lambda A$ for every sample and
$\lambda\ge1/n$. Splitting the two terms in the root yields
\[
 \lambda_S\le C\left[
 \theta_n^{1/v}\widehat R_S(g)^{1-\rho/v}
 +\theta_n^{1/\rho}(2A+b_n)^{v/\rho-1}\right].
\]
Consequently, with $\bar A=A+b_n+B_0+1$, the empirical oracle is
\[
 \widehat R_S(g)+C\left[
 \bar A\widehat R_S(g)^{1-\rho/v}\theta_n^{1/v}
 +\bar A^{v/\rho}\theta_n^{1/\rho}\right]+n^{-3}.
\]
If $\|g\|_\infty^2\le C A^r$, Bernstein's scalar
teacher concentration is absorbed by those same terms, after including
confidence logarithms in $\bar A$. To check leading coefficient one,
put $D=\bar A^r\theta_n^{1/\rho}$. Bernstein gives
$\widehat R(g)-R(g)\le C\sqrt{R(g)D}+CD$, while
$\sqrt{RD}\le R^{1-1/r}D^{1/r}+D$ for $r\ge2$.
Also $\widehat R(g)\le2R(g)+CD$, which handles the sublinear power
without multiplying the leading population-risk term. Thus
\begin{equation}
 R(\mathbb E_\Pi z)\le R(g)+C\left[
 \bar A R(g)^{1-1/r}\theta_n^{1/v}
 +\bar A^r\theta_n^{1/\rho}\right]+n^{-3}.
 \label{eq:localized-population-oracle}
\end{equation}
The fallback is absorbed by the second remainder when its branch
condition fails, after enlarging constants. This theorem is uniform in
all eligible widths/scales; the comparison teacher's scalar
concentration event is common to them.


\section{Applying the localized oracle to the Gaussian network}
\label{sec:assembly}
\begin{proof}[Proof of Theorem~\ref{thm:main}]
At $H=1$, the localized capacity argument in Section~\ref{sec:one-hidden} proves the required raw-risk certificate with $(v,\rho)=(2,1)$; its teacher conversion below is unchanged. Assume $H\ge2$ for the program construction. The centered independent readout and bounded hidden activations give
$\E_\mu\widehat R\le2$ for every sample, width, and scale. Take $M=2$ in Section~\ref{sec:localized-direct-oracle}.

Apply Proposition~\ref{prop:programs} with the proof measure
$\xi_S=(D_X+\widehat D_{X,S})/2$. Quantize its fixed compact program family in uniform norm at accuracy $c_Ha$. After rescaling $a$ by a depth-dependent constant, this gives exactly the finite dictionaries required there, with $b=1$, $p=H^2+2H-3$, $q=4H-2$, and a dimension factor polynomial in $d+1$. The logarithmic factor in \eqref{eq:program-entropy} is absorbed by the larger explicitly chosen $\theta_n$. Thus
\[
 \rho=(q+2)/2=2H,\qquad v=\max\{p+1,2\rho\}=v_H.
\]
The same dictionaries cover every finite width and Gaussian scale. The statistical event is consequently common to all these learners; there is no union over a continuum of priors.

For a fixed teacher, Lemma~\ref{lem:teacher} and the variational formula for free energy give, for $\lambda\ge1/n$,
\[
 F_S(\lambda)\le\widehat R_S(g)+\lambda[1+2A_g(\lambda)]
 \le\widehat R_S(g)+\lambda[1+2A_g(1/n)].
\]
Also $\|g\|_\infty^2\le m^2B^2\le E_g$, so its output envelope is controlled by the comparison resource. The scalar Bernstein event for this one teacher is common to all learner widths and scales. Allocate half the failure probability to it and half to the fixed-dictionary event. The population-oracle argument of \eqref{eq:localized-population-oracle}, with the enlarged resource $\mathcal A_{g,n}$, proves \eqref{eq:main-oracle}. The fallback branch is covered as in that argument.

Every finite posterior has finite readout moments because its density is bounded by $Z^{-1}$ relative to its Gaussian prior. Jensen's inequality therefore proves the same bound for the ordinary raw posterior mean.
\end{proof}

\begin{proof}[Proof of Corollary~\ref{cor:rates}]
At fixed resources, $\mathcal A_{g,n}=O_H(\log(n/\delta))$ and $\theta_n$ is $n^{-1}$ times logarithmic factors. This proves the two rates. For $m\le n^\alpha$ and $B\le n^\beta$, at fixed $d,H$, the resource is at most $O_H(n^{2(\alpha+\beta)}\log n)$. The first remainder vanishes when $2(\alpha+\beta)<1/v_H$; the second gives the same inequality because $r_H=v_H/\rho_H$.

A minimum-risk teacher exists in each compact bounded-weight class by dominated convergence. The zero-readout network belongs to that class, so its optimal risk is at most one. Apply the theorem to this deterministic population minimizer. The calibration does not depend on the teacher class. A countable family of comparison classes may be covered by summable confidence allocations, adding only the corresponding logarithmic confidence costs.
\end{proof}

\paragraph{Measurability.}
Finite program families are compact continuous images of bounded parameter boxes and have countable uniformly dense subsets. Actual-network norm tests and residual tests are Borel in the finite weights and sample; spectral bases can be chosen measurably from finite Gram matrices. Population integrals are measurable by bounded convergence. The statistical suprema are over fixed dictionaries, whereas sample-dependent good sets enter only prior integrals whose bounds hold for every fixed sample. No proof-program selection is part of the learner.

\section{One hidden layer}
\label{sec:one-hidden}
One hidden layer admits a sharper elementary capacity estimate for the same activation class, including nonzero $\phi(0)$.

\begin{lemma}[Normalized ridge atoms]
\label{lem:normalized-atoms}
Put $R_d=\sqrt{d+1}$. The class
\[
 \mathcal A=\left\{x\mapsto
   \frac{\phi(w\cdot\bar x)-\phi(0)}{\|w\|_2}:w\ne0\right\}
\]
has envelope $R_d$ and worst-case empirical Rademacher complexity at most $C_\phi(d+1)/\sqrt n$.
\end{lemma}
\begin{proof}
Write $w=ru$ with $\|u\|=1$ and $r>0$. Extend the atom continuously to $r=0$ by $\phi'(0)u\cdot\bar x$. The integral formula
\[
 A_{r,u}(x)=\int_0^1\phi'(tr(u\cdot\bar x))(u\cdot\bar x)dt
\]
shows that it is $R_d$-Lipschitz in $u$ and $\tfrac12\sup|\phi''|R_d^2$-Lipschitz in $r$, uniformly over inputs. The first claim also follows directly from the Lipschitz property of $\phi$ before dividing by $r$. Since $|\phi|\le1$, every $r>2/a$ gives an atom uniformly bounded by $a$.

For $0<a\le R_d$, cover the unit sphere in $u$ at scale $a/(4R_d)$ and the interval $0\le r\le2/a$ at scale $a/[2(1+\sup|\phi''|)R_d^2]$. Together with the zero function for the tail, this gives
\[
 \log\mathcal N(Ca,\mathcal A,\|\cdot\|_\infty)
 \le C(d+1)\log(C_\phi R_d/a).
\]
Dudley's entropy integral is at most $C_\phi R_d\sqrt{(d+1)/n}$, proving the assertion.
\end{proof}

For a one-hidden-layer network on the norm event, write $c_i=Nv_i$, and use finite averages over its $N$ rows. Then
\[
 z=\phi(0)\E c+\E[c\|w\|A_w],\qquad
 |\E c|\le\sqrt K,\qquad
 \E|c|\|w\|\le\sqrt{\E c^2\E\|w\|^2}\le K.
\]
Thus the output class is contained in a constant-function interval of radius $\sqrt K$ plus $K$ times the symmetric convex hull of $\mathcal A$. Its Rademacher complexity is at most $C_\phi(d+1)K/\sqrt n$, while actual outputs are bounded by $\sqrt K$.

\paragraph{The localized shallow certificate.}
The same norm-good sets have prior complements at most $C_\phi e^{-c_\phi\omega K}$ at every width and scale. On the objective region $\widehat R+\beta\lambda K\le u$, the smooth-square-loss bound \eqref{eq:smooth-loss} below gives deviation at most $\beta\lambda K$ whenever
\[
 n\ge C_\phi(d+1)^2\log^c(C_\phi n/\delta)\,u/\lambda^2.
\]
Its leading term is $C_\phi\sqrt u\,K/\sqrt n$; the squared term is absorbed using $\lambda K\le C u$. Confidence terms are smaller because $K$ has a positive constant lower cutoff. Prebuild dyadic norm and risk levels. The localized prior-tail and entropy argument in Section~\ref{sec:localized-direct-oracle} then gives $\int R\,d\Pi\le F_S(\lambda)+C_\phi\lambda+n^{-3}$ at the root $\lambda^2=\theta_n[2F_S(\lambda)+b_n\lambda]$. This proves the main theorem's $H=1$ case, with $(v_1,\rho_1)=(2,1)$ and oracle remainder
\[
 C_\phi\big[\mathcal A_{g,n}\sqrt{R(g)\theta_n}+\mathcal A_{g,n}^2\theta_n\big]+n^{-3}.
\]
In particular its zero-risk rate is $\widetilde O(n^{-1})$ for the full bounded smooth activation class. The fixed-calibration monotone result below remains useful because it has a larger growing-comparison range.

We record the resulting fixed-calibration statement separately from the evidence-calibrated deeper theorem. For $0<\lambda\le\epsilon\le1$ define
\[
 \begin{aligned}
 \ell_n&=\log\frac{en(d+1)}{\lambda\epsilon\delta},\\
 \mathcal B_g&=R(g)+2\lambda A_g(\lambda)
 +C_\phi[\epsilon+\lambda\log(C_\phi/\delta)]\\
 &\quad +(U_g+1)\sqrt{\frac{2R(g)\log(8/\delta)}n}
 +\frac{2(U_g+1)^2\log(8/\delta)}{3n}.
 \end{aligned}
\]
\begin{proposition}[Fixed-calibration shallow oracle]
\label{prop:shallow}
For sufficiently small $\lambda$, the posterior with likelihood variance $n\lambda\tau^2/N$ satisfies
$\int R(z_w)d\Pi\le\mathcal B_g$, and the same bound holds for the raw posterior mean, on an event of probability at least $1-\delta$ common to every $N\ge m$ and $0<\tau\le1$, provided
\[
 n\ge C_\phi(d+1)^2\ell_n^{18}/\lambda^2.
\]
If the activation is also monotone, the smaller sufficient budget is
\[
 n\ge C_\phi(d+1)\ell_n^{18}/(\lambda\epsilon).
\]
Consequently the generic fixed-teacher rate is $\widetilde O(n^{-1/2})$. For monotone activations, comparison width $n^\alpha$ and coefficient cap $n^\beta$ are allowed with vanishing excess whenever $\alpha+\beta<1/2$.
\end{proposition}
\begin{proof}
We give the statistical argument to specify the calibration and its uniformity. A standard smooth-square-loss Rademacher bound \cite{sst} says that a fixed function class of envelope $\sqrt K$ and worst-case empirical Rademacher complexity $\mathfrak r_n$ satisfies, on one sample event,
\begin{equation}
 \begin{aligned}
 R(f)-\widehat R(f)\le C\bigg[&\sqrt{\widehat R(f)}
 \left(\log^{3/2}(2n)\mathfrak r_n+\sqrt{(K+1)\log(2/\delta')/n}\right)\\
 &+\log^3(2n)\mathfrak r_n^2+(K+1)\log(2/\delta')/n\bigg].
 \end{aligned}
 \label{eq:smooth-loss}
\end{equation}
Apply it at dyadic norm levels $K$ up to a sufficiently large multiple of $1/\lambda$. The preceding normalized-atom estimate gives $\mathfrak r_n\le C_\phi(d+1)K/\sqrt n$. The generic budget makes the deviation, for empirical risk at most eight, at most $\lambda K/64+C_\phi\epsilon$.

For a monotone activation a sharper dictionary bound applies. If its endpoint limits are $l_-,l_+$, it is the constant $(l_++l_-)/2$ plus a positive mixture of shifted signs of total mass $(l_+-l_-)/2$. Their absolute masses sum to at most one. A decreasing activation is handled by changing sign. Thus its ridge dictionary lies in the symmetric convex hull of affine-halfspace dichotomies. Sauer--Shelah and the finite-class Rademacher bound give dictionary complexity $C\sqrt{(d+1)\log(n+1)/n}$. The readout has $\ell^1$ norm at most $\sqrt K$, giving output-class complexity $C\sqrt{K(d+1)\log(n+1)/n}$. The identity
\[
 \sup_{K>0}\frac{K}{(\epsilon+\lambda K)^2}=\frac1{4\lambda\epsilon}
\]
then yields the monotone budget and the same deviation bound.

For completeness, pass from these level bounds to the posterior without restricting the prior. Norm-good sets have complements at most $C e^{-\omega K/32}$ at every width. Define $Q_S$ to be twice the largest failed level, with a constant lower cutoff, and use the largest level as a terminal exceptional state. Prior tail integration gives
$\log\E_\mu e^{\omega Q_S/128}\le C_\phi\omega$.
Entropy duality and the teacher comparison imply
\[
 \int[\widehat R+\lambda Q_S/64]d\Pi
 \le\widehat R(g)+2\lambda A_g(\lambda)+C_\phi\lambda.
\]
The centered readout gives $Z\ge e^{-\omega/\lambda}$, while $R(z)\le2N\|v\|^2+2$ and the Gaussian readout imply $\E_\mu R^2\le24$. Therefore the posterior risk on empirical risk above eight, and on the terminal norm-bad event, is at most $C_\phi e^{-c_\phi/\lambda}$. Integrate the level deviation on the remaining states and apply the scalar teacher Bernstein inequality. This proves the proposition on one sample event for all widths/scales. For the monotone growing-resource claim choose $\lambda=n^{-t}$ and $\epsilon=n^{-u}$ with $2(\alpha+\beta)<t<1$ and $0<u\le t$, $t+u<1$.
\end{proof}

Tanh and sigmoid satisfy the monotone hypothesis. Bounded smooth nonmonotone activations such as sine are covered by the generic atom argument; their unrestricted ridge dictionaries need not have finite VC dimension.

\paragraph{A second fixed-calibration option at two hidden layers.}
The deeper program argument also permits a useful alternative calibration when $H=2$. For $0<\lambda\le\epsilon\le1$, the same oracle expression $\mathcal B_g$, with depth-dependent constants, follows if
\[
 n\ge C_H(d+1)^2\ell_n^{18}\lambda^{-6}\epsilon^{-2}.
\]
Indeed use finite dictionaries from Proposition~\ref{prop:programs} at approximation accuracy $a=c_H(\epsilon+\lambda K)$ and apply finite-class Bernstein at each dyadic norm level. Since $p_H+1=6$ and $q_H+2=8$, its sample budget is bounded by
\[
 C_H(d+1)^2\ell_n^{18}\sup_{K>0}
 \frac{K^6}{(\epsilon+\lambda K)^8}
 \le C_H(d+1)^2\ell_n^{18}\lambda^{-6}\epsilon^{-2}.
\]
The root-risk approximation transfer costs $O(a)$ on bounded empirical-risk levels. The same prior-tail and variational argument just proved applies to the sample-dependent good sets from Proposition~\ref{prop:programs}, whose mass bounds hold for each fixed sample. It therefore covers all finite widths and scales. This option gives the larger growing-resource range $\alpha+\beta<1/12$, by taking $6t+2u<1$ and $2(\alpha+\beta)<t$. The evidence-calibrated main theorem instead supplies the sharper $n^{-1/4}$ zero-risk rate without declaring a comparison class.

\section{Sharper tanh guarantees above mean-field scaling}
\label{sec:above-one}
The all-scale theorem also covers mean field. For tanh, stronger scalings permit an additional localization of a spectral-truncation argument. This section improves the risk dependence of the known finite-width spectral/small-width comparison; it does not claim that a rate denominator linear in depth for fixed $\gamma>1$ is new.

Fix $\phi=\tanh$, $H\ge2$, and $\gamma>1$, with $\tau^2=N^{1-\gamma}$. Put
\begin{equation}
 h=\max\{H,3\},\qquad
 v_\gamma=3+\frac{\gamma(h-1)}{\gamma-1},\qquad
 \rho_\gamma=1+\frac\gamma{\gamma-1},\qquad r_\gamma=v_\gamma/\rho_\gamma.
 \label{eq:gamma-powers}
\end{equation}
Here $r_\gamma\ge2$. Use the evidence-root rule \eqref{eq:main-calibration} with $(v_H,\rho_H)$ replaced by $(v_\gamma,\rho_\gamma)$ and sufficiently large constants depending on $H,\gamma$. The prior is unchanged.

\begin{theorem}[Optimistic spectral comparison]
\label{thm:above-one}
For every fixed teacher $g$, with probability at least $1-\delta$, simultaneously for every $N\ge m$, the raw posterior risk and the raw posterior mean risk are at most
\begin{equation}
 R(g)+C_{H,\gamma}\left[
 \overline A_{g,n}R(g)^{1-\rho_\gamma/v_\gamma}
           \theta_n^{1/v_\gamma}
 +\overline A_{g,n}^{v_\gamma/\rho_\gamma}
           \theta_n^{1/\rho_\gamma}\right]+n^{-3},
 \label{eq:above-one-oracle}
\end{equation}
where $\overline A_{g,n}=C_{H,\gamma}+2A_g(1/n)+\log(C_{H,\gamma}n/\delta)$ and $\theta_n$ is $n^{-1}$ times the polynomial dimension and logarithmic factors from the main theorem, with enlarged constants. For $\gamma\ge2$ the same statement holds with $(v,\rho)=(H,1)$.
\end{theorem}

For example, at $\gamma=3/2$ and $H\ge3$, $(v_\gamma,\rho_\gamma)=(3H,4)$. At fixed resources the zero-risk rate is then $n^{-1/4}$ up to logarithms, independently of depth. Its comparison-resource exponent is $3H/4$, and the constants still depend on depth. Thus this does not yield a polynomial bound jointly in comparison width and depth.

\subsection{A localized finite-width PAC--Bayes certificate}
Write $F=F_S(\lambda)$ and $u=2F+b_n\lambda$, where
$b_n=C_{H,\gamma}\log(C_{H,\gamma}n/\delta)$, and work on the ordinary root branch $1/n\le\lambda\le1/b_n$, $u\le5$.
For each integer width, prebuild inverse-temperature values $\beta_j=(\omega/2)2^j$ through $\omega n$ and squared clipping radii $M_k^2=2^k$ through $C_H n$. Allocate confidence proportional to $\delta/N^2$ and uniformly over these $O(\log^2 n)$ grid pairs.

For the selected $\lambda$, choose $\beta_*\in[\omega/(2\lambda),\omega/\lambda]$ and a grid radius with $C_Hu/\lambda\le M^2\le2C_Hu/\lambda$. These choices are within the grids since $u\ge b_n\lambda$. The proof-clipped squared loss has range $b=(M+1)^2\le C_Hu/\lambda$.

For any bounded loss $\ell\in[0,b]$ and $s=\beta_*b/n$, convexity gives
\[
 \E e^{-s\ell/b}\le\exp[-(1-e^{-s})\E\ell/b].
\]
Independence, integration over the fixed Gaussian prior, Markov's inequality, and entropy duality give, on the grid event, every law $Q$ the bound
\[
 \int R(z^{[M]})dQ\le\frac{s}{1-e^{-s}}
 \left[\int\widehat R(z^{[M]})dQ+
             \frac{\KL(Q\Vert\mu)+t}{\beta_*}\right],
\]
where $t\le C_H\log(C_HnN/\delta)$. For the actual posterior, the clipped empirical objective in brackets without $t$ is at most its raw free energy $F$: clipping does not increase square loss and $\beta_*$ is at least its actual inverse temperature. If $s\le1$, then $s/(1-e^{-s})\le1+s$, so
\begin{equation}
 \int R(z^{[M]})d\Pi\le F+
       \frac{C_H\omega u^2}{n\lambda^2}
       +C_H\lambda\log(C_Hn/\delta).
 \label{eq:localized-finite-width}
\end{equation}
Here $\log N/\omega\le1$ because $\omega=N^\gamma\ge N$.

For raw tails put $r=N\|v\|^2$. The Gaussian readout gives
$\mu(r>M^2)\le e^{-\omega M^2/8}$ and $\E_\mu R^2\le24$. Exact evidence, rather than a constant comparison risk, gives
\[
 \int R(z)\mathbf1_{r>M^2}d\Pi
 \le\sqrt{24}\exp[\omega F/(2\lambda)-\omega M^2/16]
 \le n^{-3}
\]
after choosing the radius constant sufficiently large. Thus whenever
\begin{equation}
 n\ge C_H\omega u^2\lambda^{-3}\log(C_Hn/\delta),
 \label{eq:small-localized-budget}
\end{equation}
the raw risk is at most $F+C_H\lambda\log(C_Hn/\delta)+n^{-3}$. The condition also ensures $s\le1$. The same event handles every eligible data-dependent calibration and every width.

\subsection{The spectral certificate at larger widths}
Retain from each internal matrix only singular values greater than $8\tau$, obtaining $T^\ell$, and leave the first layer and readout unchanged. Let $\widetilde z$ be the resulting network and set
\[
 q=K_0+\sum_{\ell=2}^H\|T^\ell\|_F^2+N\|v\|^2.
\]
For each internal Gaussian matrix, the operator-norm ball of radius $4\tau$ has probability at least one half. Its squared Frobenius distance is $\sum_j(s_j-4\tau)_+^2$, and the retained energy is at most four times this quantity. Gaussian isoperimetry and integration therefore give
\[
 \E\exp(\omega\|T^\ell\|_F^2/16)\le3/2.
\]
The readout chi-square moment gives
$\log\E e^{\omega q/16}\le C_H\omega$.
The same estimate at a smaller exponent controls raw-risk-weighted moments, because $R(z)\le2q+2$.

Normalized hidden-state perturbations and the Lipschitz property of tanh give
\begin{equation}
 \|z-\widetilde z\|_\infty\le C_H\tau q^{(H-1)/2},
 \qquad |z|,|\widetilde z|\le\sqrt q.
 \label{eq:localized-spectral-error}
\end{equation}
The union over all widths of tanh networks with internal-plus-readout energy at most $K$ and unrestricted first-layer weights has Rademacher complexity
\begin{equation}
 \mathfrak r_n\le C_HK^{H/2}\sqrt{(d+1)\log(n+1)/n}.
 \label{eq:localized-spectral-rad}
\end{equation}
For clarity, this bound uses more than bounded smoothness. First-layer tanh features lie in the convex hull of affine-halfspace signs. Their Rademacher complexity is at most $\sqrt{2(d+1)\log(n+1)/n}$. If a symmetric convex class lies in $[-1,1]$, adjoining all variable clips $\clip_{[-t,t]}$, $0\le t\le1$, increases its complexity by at most $1/(2\sqrt n)$, by contraction and the identity $\clip_{[-t,t]}u=(|u+t|-|u-t|)/2$. The identity
\[
 \frac{\tanh(su)}s=
 \int_0^1\clip_{[-t,t]}(u)\,2s\tanh(st)\operatorname{sech}^2(st)dt
 +\operatorname{sech}^2(s)\clip_{[-1,1]}(u)
\]
for $|u|\le1$ has nonnegative coefficients of total mass one. It propagates the convex-hull estimate through successive layers. The final path-weight sum is at most
$\sqrt N\|v\|\prod_{\ell=2}^H\|T^\ell\|_F\le(K/H)^{H/2}$.
This proves \eqref{eq:localized-spectral-rad}. We do not apply this unrestricted first-layer argument to nonmonotone activations.

On $J=\widehat R(z)+\beta_0\lambda q\le u$, the energy is at most $C_Hu/\lambda$. A sufficient condition for \eqref{eq:localized-spectral-error} to be at most $c_H\lambda q/\sqrt u$ throughout this region is
\begin{equation}
 \tau^2\le c_H\lambda^{h-1}u^{2-h},\qquad h=\max\{H,3\}.
 \label{eq:localized-spectral-scale}
\end{equation}
For $H\ge3$, maximize $q^{(H-3)/2}$ at the upper energy endpoint; for $H=2$, its maximum is at the constant lower cutoff. Apply the smooth-square-loss inequality \eqref{eq:smooth-loss} to fixed energy classes at dyadic levels. At risk at most $C_Hu$, its deviation is at most $\beta_0\lambda K$ if
\[
 n\ge C_H(d+1)\log^c(C_Hn/\delta)\,
                  u\lambda^{-2}K^{H-2}.
\]
Maximizing over $K\le C_Hu/\lambda$ gives the sufficient budget
\begin{equation}
 n\ge C_H(d+1)\log^c(C_Hn/\delta)\,
                    u^{H-1}\lambda^{-H}.
 \label{eq:localized-spectral-budget}
\end{equation}
The linear and confidence terms are absorbed using $\lambda K\le C_Hu$. Equation~\eqref{eq:localized-spectral-scale} transfers empirical and population root risks with cost at most $c_H\lambda K$.

The localized tail and entropy argument of Section~\ref{sec:localized-direct-oracle} now applies with complexity $q$. On $J>u$ its raw-risk contribution is at most $n^{-3}$; on $J\le u$ integrate the level deviation. It follows that the raw risk is at most $F+C_H\lambda+n^{-3}$. All statistical classes are independent of width and scale. A grid of risk and calibration values supplies one event for all choices made from the sample.

\subsection{Combining the two certificates}
Condition \eqref{eq:localized-spectral-scale} holds once
\[
 N\ge C_H[u^{h-2}\lambda^{-(h-1)}]^{1/(\gamma-1)}.
\]
Below this cutoff, $\omega=N^\gamma$ is at most the $\gamma$th power of its right side, so \eqref{eq:small-localized-budget} follows from
\[
 n\ge C_{H,\gamma}(d+1)\log^c(C_{H,\gamma}n/\delta)\,
                   u^{v_\gamma-\rho_\gamma}\lambda^{-v_\gamma}.
\]
This also dominates \eqref{eq:localized-spectral-budget}: their ratio is
$u^{\rho_\gamma-1}(\lambda/u)^{v_\gamma-H}$, which is bounded since $u\le5$, $\lambda/u\le1/b_n$, and $v_\gamma\ge H$.
The evidence root enforces this budget. Apply the teacher comparison and the same coefficient-one Bernstein conversion as in Section~\ref{sec:localized-direct-oracle}, absorbing the logarithmic certificate coefficient into $\overline A_{g,n}$. This proves \eqref{eq:above-one-oracle}.

When $\gamma\ge2$, use the full internal Frobenius energy in $q$ instead. Its Gaussian exponential moment costs $O_H(N^2)$ in the logarithm; divided by $\omega=N^\gamma$, this is bounded uniformly. No spectral approximation is needed. The only sample condition is \eqref{eq:localized-spectral-budget}, giving $(v,\rho)=(H,1)$. The same teacher comparison completes the proof.

\subsection{Selecting a calibration while keeping one prior}
All branches give observable certificates of the form $F_S(\lambda)+B_n\lambda+n^{-3}$, with one sufficiently large common $B_n$, on a joint event obtained by a finite confidence split. Their fallback branches instead give their stated constant-risk certificate. Compute the applicable calibrations and select the one with smallest certificate. This still uses one architecture and its original Gaussian prior; only its likelihood variance is calibrated. On ordinary branches the certificate is nondecreasing in $\lambda$, since $F_S(\lambda)$ is nondecreasing. Thus selecting the smallest ordinary-root $\lambda$ gives the best of those certificates. Including fallback branches requires comparing their actual constant certificates as well.

The available exponent pairs for tanh, each yielding the oracle form
$A R(g)^{1-\rho/v}n^{-1/v}+A^{v/\rho}n^{-1/\rho}$ up to logarithms, are:
\begin{center}
\begin{tabular}{lll}
\toprule
Regime & $v$ & $\rho$\\
\midrule
$H=1$, every $\gamma\ge1$ & $2$ & $1$\\
$H\ge2$, every $\gamma\ge1$ & $\max\{4H,H^2+2H-2\}$ & $2H$\\
$H\ge2$, fixed $\gamma>1$ & $3+\gamma(\max\{H,3\}-1)/(\gamma-1)$ & $1+\gamma/(\gamma-1)$\\
$H\ge2$, $\gamma\ge2$ & $H$ & $1$\\
\bottomrule
\end{tabular}
\end{center}
The second row covers the full bounded smooth activation class, as does the first. The sharper last two rows here use tanh. Depending on comparison risk and resources, different pairs can be preferable; selecting a valid certificate gives the corresponding best risk bound without changing the prior.


\section{Scope and directions for improvement}
The theorem combines an approximation term with a complexity penalty under one fixed dense Gaussian prior. It allows arbitrary bounded IID labels and arbitrarily large learner widths. Neither the Gaussian prior nor the real-valued prediction is clipped or restricted.

The proof-only compressed functions use Gaussian-smoothed scalar activations, scalar clipping, and finite linear combinations. They bound the statistical capacity needed by the argument; they do not replace the learner. Good parameter sets used in the proof depend on the population input law and the sample. This is legitimate because their prior mass bounds hold for every fixed sample, while statistical uniformity is applied only to the fixed family of proof programs.

The remaining quantitative loss for positive comparison risk comes from representing several spectral directions at each layer. A sharper capacity bound for these Gaussian feature constructions could improve the quadratic depth exponent. The comparison network must have bounded weights, and constants can grow quickly with depth. No computational posterior-sampling guarantee is supplied.

These are conventional mathematical proofs, not Lean formalizations. The verification directory records independent checks and the source snapshot reviewed.

\begin{thebibliography}{9}
\bibitem{vh} Dmitry Vaintrob and Kaarel H\"anni. \emph{Width-Robust Learnability in Mean-Field Bayesian Neural Networks}. arXiv:2607.05735, 2026. \url{https://arxiv.org/abs/2607.05735}.
\bibitem{sst} Nathan Srebro, Karthik Sridharan, and Ambuj Tewari. \emph{Smoothness, Low Noise and Fast Rates}. Advances in Neural Information Processing Systems 23, 2010. \url{https://www.cs.cornell.edu/~sridharan/smoothness.pdf}.
\bibitem{dudley} R. M. Dudley. \emph{The sizes of compact subsets of Hilbert space and continuity of Gaussian processes}. Journal of Functional Analysis 1, 290--330, 1967.
\bibitem{contraction} Michel Ledoux and Michel Talagrand. \emph{Probability in Banach Spaces}. Springer, 1991.
\end{thebibliography}
\end{document}
